\documentclass[11pt]{article}
\usepackage{amsmath,amssymb}
\usepackage[margin=1in]{geometry}
\begin{document}

\section*{Step 213: Bagchi Universality Carrier}

Let \(K\subset\{1/2<\Re(s)<1\}\) be compact with connected complement, and let
\(f\) be continuous on \(K\), holomorphic inside \(K\), and nonvanishing on
\(K\).  The Bagchi/Voronin residual is
\[
  \Xi_{\rm Bagchi}(K,f)
  =\inf_{\tau\in\mathbb R}\sup_{s\in K}
  |\zeta(s+i\tau)-f(s)|.
\]
The vertical-translation operator is
\[
  T_\tau F(s)=F(s+i\tau).
\]

\subsection*{Theorem}

Voronin and Bagchi universality state that for every \(\epsilon>0\), the set
of \(\tau\) satisfying
\[
  \sup_{s\in K}|\zeta(s+i\tau)-f(s)|<\epsilon
\]
has positive lower density.  Thus \(\Xi_{\rm Bagchi}(K,f)=0\) is already
proved unconditionally.

\subsection*{CRE audit}

Because the residual is closed unconditionally, it is not equivalent to RH.
It is a value-distribution/statistical theorem about zeta translates, not a
zero-location closure theorem.

\subsection*{Derivative carriers}

Joint universality and Selberg-class universality may inform statistical
questions about \(L\)-functions, but no inherited residual of the form
\[
\text{closure} \Longleftrightarrow \mathrm{RH}
\]
is produced by Bagchi universality itself.

\subsection*{Verdict}

\[
\boxed{\texttt{V\_bagchi\_outside\_dichotomy}.}
\]

The classification is analogous to the RMT carrier from Step 188: outside the
RH-closure Dichotomy scope, not a coverage refuter.

\end{document}
